% GENERATED FILE. Edit canonical lessons, then run npm run generate:books.
% canonical-source-hash: fnv1a64:6fa3210b648a442f
% canonical-lessons: ZH-C101-chufang, ZH-C101-baicai, ZH-W101-cai, ZH-C101-kuaizi, ZH-C101-panzi, ZH-C101-dao

\chapter{In the Kitchen}
\label{ch:zh-in-the-kitchen}
\hlchaptermodality{101}

\begin{chapteropening}
\textbf{By the end of this chapter:} \emph{I can name the kitchen, cabbage, chopsticks, plates and a knife.}

Complete the last lesson of chapter 101: I can name the kitchen, cabbage, chopsticks, plates and a knife.
\end{chapteropening}

\section[chúfáng]{chúfáng — kitchen}

\label{lesson:ZH-C101-chufang}



\begin{quote}
Before the new one: say the Mandarin for library, then the Mandarin for post office.
\end{quote}

\subsection*{You'll want to know: chúfáng}
\emph{chúfáng} — \textquotedblleft{}kitchen\textquotedblright{}.

\begin{quote}
\emph{māma zài chúfáng} — Mum's in the kitchen.
\end{quote}

Its characters come later; for now it is a word for the ear.

\begin{cousinweb}
\emph{chúfáng} is built from chú and fáng 'room', both inherited from Middle Chinese, \emph{drju} and \emph{bjang} (Guangyun, in Baxter's transcription); the dictionary glosses the first as the house where food is cooked.
\end{cousinweb}

\begin{grammarlens}[title={What you've built}]
One more word for everyday things.
\end{grammarlens}

\subsection*{Your turn}
\begin{itemize}
\raggedright
  \item \emph{Say it:} \emph{chúfáng}
  \item \emph{Say it:} \emph{chúfáng}, once more
  \item \emph{Say it:} \emph{māma zài chúfáng} — Mum's in the kitchen.
  \item \emph{Recall:} say the Mandarin for library, then the Mandarin for post office, then say \emph{chúfáng} again
\end{itemize}

\begin{tcolorbox}[breakable,colback=teal!4,colframe=teal!35!black,title={Before you move on}]
What does chúfáng mean? (\textquotedblleft{}Kitchen\textquotedblright{}.) Say it once more.
\end{tcolorbox}
\section[\texorpdfstring{\zh{白菜} (báicài)}{白菜 (báicài)}]{\zh{白菜} (báicài) — Chinese cabbage}

\label{lesson:ZH-C101-baicai}



\begin{quote}
Before the new one: say the Mandarin for post office, then the Mandarin for kitchen.
\end{quote}

\subsection*{You'll want to know: \zh{白菜}}
\textbf{\zh{白菜}} — \emph{báicài} — \textquotedblleft{}Chinese cabbage\textquotedblright{}.

\begin{quote}
\textbf{\zh{我去洗白菜}} \emph{wǒ qù xǐ báicài} — I'll go and wash the cabbage.
\end{quote}

The new character, in pieces: \zh{菜} — the grass piece on top, three strokes; under it a hand reaching down, four strokes; and at the foot the tree piece of \zh{机}, standing alone. The next lesson puts it in your hand.

\begin{cousinweb}
\textbf{\zh{白菜}} (\emph{báicài}) is built from \zh{白} 'white' and \zh{菜} 'vegetable', both inherited from Middle Chinese, \emph{baek} and \emph{tshojH} (Guangyun, in Baxter's transcription). The dictionary says every plant that can be eaten is called \emph{tshojH}; the white one is probably named for its pale stalks.
\end{cousinweb}

\begin{grammarlens}[title={What you've built}]
One more word for everyday things.
\end{grammarlens}

\subsection*{Your turn}
\begin{itemize}
\raggedright
  \item \emph{Say it:} \emph{báicài}
  \item \emph{Say it:} \emph{báicài}, once more
  \item \emph{Say it:} \emph{wǒ qù xǐ báicài} — I'll go and wash the cabbage.
  \item \emph{Recall:} say the Mandarin for post office, then the Mandarin for kitchen, then say \emph{báicài} again
\end{itemize}

\begin{tcolorbox}[breakable,colback=teal!4,colframe=teal!35!black,title={Before you move on}]
What does \zh{白菜} mean? (\textquotedblleft{}Chinese cabbage\textquotedblright{}.) Say it once more.
\end{tcolorbox}
\section[\texorpdfstring{\zh{菜} (cài)}{菜 (cài)}]{\zh{菜} — the grass piece, a hand, the tree piece}

\label{lesson:ZH-W101-cai}



\begin{quote}
The word first, then the character inside it.

\begin{itemize}
\raggedright
  \item \emph{Recall:} say \emph{báicài} — \textbf{R1}, one lesson back, before it decays
\end{itemize}

\begin{quote}
\zh{机} \zh{白}
\end{quote}

\emph{Read it:} each one aloud, with nothing said first — \emph{jī}, \emph{bái}
\end{quote}

\subsection*{Script — the grass piece, a hand, the tree piece}
\hlblockfigure{\detokenize{figures/ZH-W101-cai-filmstrip.pdf}}{How \zh{菜} is written, stroke by stroke}

\begin{quote}
\zh{菜}
\end{quote}

\textbf{Three pieces, top to bottom.} Count the pieces, not the strokes.

\textbf{The grass piece, three strokes}: a horizontal across the top, then two short verticals through it, left before right.

\textbf{The hand, four strokes}: a left-falling stroke across, two dots under it, left then middle, and a short left-falling stroke on the right.

\textbf{The tree piece, four strokes}, full width at the foot: a horizontal, a vertical down through it, a left-falling stroke and a long right-falling one.

3 + 4 + 4 = eleven.

\subsection*{Your turn}
\begin{itemize}
\raggedright
  \item \emph{Write it:} \zh{菜}, all eleven, with the model in front of you
  \item \emph{Cover:} write it again from memory, and count the pieces, not the strokes
\end{itemize}

\begin{tcolorbox}[breakable,colback=teal!4,colframe=teal!35!black,title={Before you move on}]
How many pieces in \zh{菜}? (\textbf{Three} — the grass piece, a hand, the tree piece.) Which character you know has the tree piece? (\textbf{\zh{机}}, on its left.) How does \zh{菜} end? (\textbf{With a long right-falling stroke}.)

Source: \href{https://raw.githubusercontent.com/chanind/hanzi-writer-data/68d10a4b21150cae5e1ebbd223eed289cf32d90c/data/\%E8\%8F\%9C.json}{Hanzi Writer Data \zh{菜}, pinned stroke paths and medians}.
\end{tcolorbox}
\section[kuàizi]{kuàizi — chopsticks}

\label{lesson:ZH-C101-kuaizi}



\begin{quote}
Before the new one: say the Mandarin for kitchen, then the Mandarin for cabbage.
\end{quote}

\subsection*{You'll want to know: kuàizi}
\emph{kuàizi} — \textquotedblleft{}chopsticks\textquotedblright{}.

\begin{quote}
\emph{wǒ de kuàizi ne} — Where are my chopsticks?
\end{quote}

Its characters come later; for now it is a word for the ear.

\begin{cousinweb}
\emph{kuàizi} is built from kuài and the noun ending zi. The ending is inherited from Middle Chinese \emph{tsiX} (Guangyun, in Baxter's transcription); no dictionary we use records an older form of the first, so its history is not settled.
\end{cousinweb}

\begin{grammarlens}[title={What you've built}]
One more word for everyday things.
\end{grammarlens}

\subsection*{Your turn}
\begin{itemize}
\raggedright
  \item \emph{Say it:} \emph{kuàizi}
  \item \emph{Say it:} \emph{kuàizi}, once more
  \item \emph{Say it:} \emph{wǒ de kuàizi ne} — Where are my chopsticks?
  \item \emph{Recall:} say the Mandarin for kitchen, then the Mandarin for cabbage, then say \emph{kuàizi} again
\end{itemize}

\begin{tcolorbox}[breakable,colback=teal!4,colframe=teal!35!black,title={Before you move on}]
What does kuàizi mean? (\textquotedblleft{}Chopsticks\textquotedblright{}.) Say it once more.
\end{tcolorbox}
\section[pánzi]{pánzi — plate; dish}

\label{lesson:ZH-C101-panzi}



\begin{quote}
Before the new one: say the Mandarin for cabbage, then the Mandarin for chopsticks.
\end{quote}

\subsection*{You'll want to know: pánzi}
\emph{pánzi} — \textquotedblleft{}plate; dish\textquotedblright{}.

\begin{quote}
\emph{zhège pánzi hěn gānjìng} — This plate is nice and clean.
\end{quote}

Its characters come later; for now it is a word for the ear.

\begin{cousinweb}
\emph{pánzi} is built from pán and the noun ending zi, both inherited from Middle Chinese, \emph{ban} and \emph{tsiX} (Guangyun, in Baxter's transcription); the dictionary glosses the first as the name of a vessel.
\end{cousinweb}

\begin{grammarlens}[title={What you've built}]
One more word for everyday things.
\end{grammarlens}

\subsection*{Your turn}
\begin{itemize}
\raggedright
  \item \emph{Say it:} \emph{pánzi}
  \item \emph{Say it:} \emph{pánzi}, once more
  \item \emph{Say it:} \emph{zhège pánzi hěn gānjìng} — This plate is nice and clean.
  \item \emph{Recall:} say the Mandarin for cabbage, then the Mandarin for chopsticks, then say \emph{pánzi} again
\end{itemize}

\begin{tcolorbox}[breakable,colback=teal!4,colframe=teal!35!black,title={Before you move on}]
What does pánzi mean? (\textquotedblleft{}Plate; dish\textquotedblright{}.) Say it once more.
\end{tcolorbox}
\section[dāo]{dāo — knife}

\label{lesson:ZH-C101-dao}



\begin{quote}
Before the new one: say the Mandarin for chopsticks, then the Mandarin for plate.
\end{quote}

\subsection*{You'll want to know: dāo}
\emph{dāo} — \textquotedblleft{}knife\textquotedblright{}.

\begin{quote}
\emph{dāo zài zhuōzi shàngbian} — The knife is on the table.
\end{quote}

Its characters come later; for now it is a word for the ear.

\begin{cousinweb}
\emph{dāo} is inherited from Middle Chinese \emph{taw} (Guangyun, in Baxter's transcription), which the dictionary glosses as a weapon, a blade for cutting things down. In the kitchen it is the knife.
\end{cousinweb}

\begin{grammarlens}[title={What you've built}]
One more word for everyday things.
\end{grammarlens}

\subsection*{Your turn}
\begin{itemize}
\raggedright
  \item \emph{Say it:} \emph{dāo}
  \item \emph{Say it:} \emph{dāo}, once more
  \item \emph{Say it:} \emph{dāo zài zhuōzi shàngbian} — The knife is on the table.
  \item \emph{Recall:} say the Mandarin for chopsticks, then the Mandarin for plate, then say \emph{dāo} again
\end{itemize}

\begin{tcolorbox}[breakable,colback=teal!4,colframe=teal!35!black,title={Before you move on}]
What does dāo mean? (\textquotedblleft{}Knife\textquotedblright{}.) Say it once more.
\end{tcolorbox}
